\chapter{Induction}\label{ch:induction}
\index{induction}

A claim about every $n \ge 1$ cannot be proved one $n$ at a time --- there are
infinitely many. Induction is the machine that handles all of them in one
argument: prove the first case by hand, then prove that each case carries the
next. After this chapter you will be able to run that machine in its two
standard gears (ordinary and strong induction), recognise its disguise
(the minimal-counterexample argument), audit other people's uses of it ---
and construct a geometric induction like the tromino tiling, not just an
algebraic one. Chapters 5, 7 and 9 run on it constantly.

\section{The domino idea and the template}\label{sec:domino}

Suppose $P(n)$ is a statement about an integer $n$. To establish
``$P(n)$ for all $n \ge b$'' by induction, you prove exactly two things:

\begin{definition}[Principle of mathematical induction]\index{induction!principle of}
If (i) \emph{base case}: $P(b)$ is true, and (ii) \emph{inductive step}: for
every $n \ge b$, $P(n) \Rightarrow P(n+1)$, then $P(n)$ holds for all $n \ge b$.
Inside the step, the assumed $P(n)$ is called the \emph{inductive
hypothesis}\index{inductive hypothesis} (I.H.).
\end{definition}

\begin{figure}[ht]\centering
% Figure: the induction machine --- base case plus a step that never fails.
% Source for Figure~\ref{fig:dominoes}. TikZ only; regenerable by compiling main.tex.
\begin{tikzpicture}[
  dom/.style={draw,rounded corners=1.5pt,fill=graybg,minimum width=4mm,minimum height=9mm,inner sep=0pt},
  fell/.style={draw,rounded corners=1.5pt,fill=accentsoft,minimum width=4mm,minimum height=9mm,inner sep=0pt,rotate=-28},
  arr/.style={->,accent,thick},
  lab/.style={font=\footnotesize\sffamily,text=slate},
]
% first domino already falling, the rest standing
\node[fell] (d1) at (0,0) {};
\node[dom] (d2) at (1.1,0) {};
\node[dom] (d3) at (2.2,0) {};
\node[dom] (d4) at (3.3,0) {};
\node[dom] (d5) at (4.4,0) {};
\node[dom] (d6) at (5.5,0) {};
\node[font=\large,text=slate] at (6.6,0.15) {\dots};
% labels
\node[lab,anchor=north] at (-0.1,-0.62) {$P(1)$};
\node[lab,anchor=north] at (1.1,-0.55) {$P(2)$};
\node[lab,anchor=north] at (2.2,-0.55) {$P(3)$};
\node[lab,anchor=north] at (3.3,-0.55) {$P(4)$};
\node[lab,anchor=north] at (4.4,-0.55) {$P(5)$};
\node[lab,anchor=north] at (5.5,-0.55) {$P(6)$};
\node[lab,anchor=east] at (-0.5,-0.62) {base:};
\draw[arr] (0.45,0.62) .. controls (0.75,0.8) .. (1.0,0.62);
\draw[arr] (1.55,0.62) .. controls (1.85,0.8) .. (2.1,0.62);
\draw[arr] (2.65,0.62) .. controls (2.95,0.8) .. (3.2,0.62);
\draw[arr] (3.75,0.62) .. controls (4.05,0.8) .. (4.3,0.62);
\draw[arr] (4.85,0.62) .. controls (5.15,0.8) .. (5.4,0.62);
\node[lab,anchor=south] at (2.9,0.95) {step: $P(n) \Rightarrow P(n+1)$, every $n$};
\node[lab,anchor=south west] at (4.6,-0.6) {};
\end{tikzpicture}

\caption{The machine's two parts. Notice that neither part alone does anything:
the fallen first domino without working links knocks over nothing, and working
links without a first domino have nothing to transmit.}
\label{fig:dominoes}
\end{figure}

The template below is compulsory in this book, on exams, and in the portfolio:
graders grade the parts, and the parts are where induction proofs break.

\begin{rulebox}[The induction template]
\index{induction!template}
\begin{enumerate}
\item \textbf{State $P(n)$} --- write the claim as a named predicate on $n$.
\item \textbf{Base case} --- prove $P(b)$ directly.
\item \textbf{Inductive step} --- ``Let $n \ge b$. Assume $P(n)$.'' Derive
$P(n+1)$, marking the exact line where the I.H. was used.
\item \textbf{Conclusion} --- ``By induction, $P(n)$ holds for all $n \ge b$.''
\end{enumerate}
\end{rulebox}

Why does the template work? Informally: $P(b)$ starts the chain and the step
transmits it; formally it is equivalent to the well-ordering principle of
Section~\ref{sec:wop}. Until then, treat the machinery as a licence granted by
Chapter~\ref{ch:proof}'s rules --- the step is an ordinary conditional proof
($P(n) \Rightarrow P(n+1)$) nested inside a bigger claim.

\section{Sums and formulas by induction}\label{sec:suminduct}

Sum formulas are induction's home turf, and they teach the single most-used
move: \emph{peel the last term} so the I.H. has something to bite.

\begin{example}\label{ex:gauss}
Prove: for every $n \ge 1$, $\displaystyle\sum_{k=1}^{n} k = \frac{n(n+1)}{2}$.
\par\smallskip
\textbf{Strategy.} Template first: name $P(n)$, prove $P(1)$, then in the step
split the sum at $k=n$ so $\sum_{k=1}^{n} k$ appears --- that subexpression
\emph{is} the I.H., and the rest is algebra.
\par\smallskip
\begin{solution}
Let $P(n)$ be the claim $\sum_{k=1}^{n} k = n(n+1)/2$.
\emph{Base.} $P(1)$: $\sum_{k=1}^{1} k = 1 = 1\cdot 2/2$. Check.
\emph{Step.} Let $n \ge 1$ and assume $P(n)$: $\sum_{k=1}^{n} k = n(n+1)/2$. Then
\[
\sum_{k=1}^{n+1} k = \left(\sum_{k=1}^{n} k\right) + (n+1)
\overset{\text{I.H.}}{=} \frac{n(n+1)}{2} + (n+1)
= \frac{(n+1)(n+2)}{2},
\]
which is exactly $P(n+1)$. \emph{Conclusion.} By induction, $P(n)$ holds for
all $n \ge 1$.
\end{solution}
\end{example}

Two things to copy: the I.H. is marked where it enters (never let a grader hunt
for it), and the goal $P(n+1)$ was kept in sight --- the whole middle of the
step is traffic between those two endpoints.

\begin{figure}[ht]\centering
% Figure: induction step on a sum -- P(n) block plus one more column.
% Source for Figure~\ref{fig:stairsum}. TikZ only; regenerable by compiling main.tex.
\begin{tikzpicture}[
  col/.style={draw,fill=rulegray!18,minimum width=6mm,anchor=south west},
  newcol/.style={draw,fill=accent!25,minimum width=6mm,anchor=south west},
  lab/.style={font=\footnotesize\sffamily,text=slate,align=center}
]
% columns of heights 1,2,3,4 (P(n) with n=4) + new column height 5
\node[col,minimum height=2mm]  at (0,0) {};
\node[col,minimum height=4mm]  at (0.65,0) {};
\node[col,minimum height=6mm]  at (1.3,0) {};
\node[col,minimum height=8mm]  at (1.95,0) {};
\node[newcol,minimum height=10mm] at (2.6,0) {};
\draw[decorate,decoration={brace,amplitude=4pt},rulegray]
  (-0.05,1.05) -- (2.5,1.05)
  node[midway,above=5pt,lab] {$P(n)$: the sum $1+2+\cdots+n$};
\node[lab] at (2.9,-0.65) {$+\,(n{+}1)$};
\node[lab] at (1.4,-1.35) {Induction step $=$ one new column added to a
sum you already control.};
\end{tikzpicture}

\caption{The induction step on a sum is one new column. Notice: the
induction hypothesis owns the whole shaded block --- the step only has to
absorb the last column.}
\label{fig:stairsum}
\end{figure}

\begin{exercise}[easy]\label{exr:sumsq}
Prove by induction: $\sum_{k=1}^{n} k^2 = \frac{n(n+1)(2n+1)}{6}$.
(Same engine: peel the last term.)
\end{exercise}

\section{Induction for inequalities}\label{sec:ineqinduct}

Induction proves inequalities too, but the step is an algebraic squeeze, not a
substitution: you chain $\le$ and $\ge$ signs from the I.H. to the target, and
each link needs its own reason.

\begin{example}\label{ex:2nvsn2}
Prove: $2^n > n^2$ for every integer $n \ge 5$.
\par\smallskip
\textbf{Strategy.} The claim is false below $5$ (check $n=2,3,4$ --- the base
window matters), so induct from $n=5$. In the step, the I.H. gives $2 \cdot
2^n > 2n^2$; the remaining job is the auxiliary inequality $2n^2 \ge (n+1)^2$,
which holds once $n \ge 3$. Every $\ge$ gets a name.
\par\smallskip
\begin{solution}
Let $P(n)$ be ``$2^n > n^2$''. \emph{Base.} $P(5)$: $32 > 25$. \emph{Step.}
Let $n \ge 5$ and assume $2^n > n^2$. Then
$2^{n+1} = 2 \cdot 2^n \overset{\text{I.H.}}{>} 2n^2$. And for $n \ge 3$,
$2n^2 - (n+1)^2 = n^2 - 2n - 1 = (n-1)^2 - 2 \ge 2 > 0$, so $2n^2 > (n+1)^2$
and $2^{n+1} > (n+1)^2$. \emph{Conclusion.} By induction, $P(n)$ holds for
all $n \ge 5$.
\end{solution}
\end{example}

\begin{tripwire}
\index{base case!wrong window}
An inequality that is true ``eventually'' is only proved from where it starts.
$2^n > n^2$ fails at $n=2,3,4$; proving $P(5)$ and the step for $n \ge 5$ is
not optional bookkeeping --- it is the claim.
\end{tripwire}

\begin{exercise}[medium]\label{exr:factorial}
Prove by induction: $n! > 2^n$ for all $n \ge 4$. (Find the smallest working
base first.)
\end{exercise}

\section{Strong induction}\label{sec:strong}
\index{strong induction}

Ordinary induction lets the step lean on $P(n)$ only. \emph{Strong induction}
lets it lean on \emph{every} earlier case:
to prove $\forall n \ge b\;P(n)$, prove base cases
$P(b), P(b{+}1), \dots, P(b{+}s{-}1)$ and the step
$\bigl(P(b) \land P(b{+}1) \land \cdots \land P(n)\bigr) \Rightarrow P(n+1)$
for all $n \ge b{+}s{-}1$. The number of base cases $s$ equals how far back the
step reaches --- Figure~\ref{fig:reachback} draws the difference.

\begin{figure}[ht]\centering
% Figure: induction vs strong induction --- what each step is allowed to lean on.
% Source for Figure~\ref{fig:reachback}. TikZ only; regenerable by compiling main.tex.
\begin{tikzpicture}[
  dot/.style={circle,draw,inner sep=1.4pt,minimum size=5.2mm,font=\footnotesize},
  bd/.style={circle,draw=accent,fill=accentsoft,inner sep=1.4pt,minimum size=5.2mm,font=\footnotesize},
  arr/.style={->,accent},
  sarr/.style={->,accent,bend left=24},
  lab/.style={font=\footnotesize\sffamily,text=slate},
]
% ---- ordinary induction ----
\node[lab,anchor=east] at (-0.3,0) {induction};
\foreach \i in {1,...,7}
  \node[dot] (a\i) at (\i*0.72,0) {\i};
\node[bd] at (0.72,0) {1};
\foreach \i [evaluate=\i as \j using int(\i+1)] in {1,...,6}
  \draw[arr] (a\i.east) -- (a\j.west);
\node[lab,anchor=west] at (5.8,0) {each step leans on the previous case only};
% ---- strong induction ----
\node[lab,anchor=east] at (-0.3,-2.1) {strong};
\foreach \i in {1,...,7}
  \node[dot] (b\i) at (\i*0.72,-2.1) {\i};
\foreach \i in {1,2,3}
  \node[bd] at (\i*0.72,-2.1) {\i};
% arrows from all earlier cases into n+1 (show into 4 and 7)
\foreach \i in {1,2,3}
  \draw[sarr] (b\i.north) to (b4.north);
\foreach \i in {4,5,6}
  \draw[sarr] (b\i.north) to (b7.north);
\node[lab,anchor=west] at (5.8,-2.1) {each step may lean on \emph{any} earlier case};
\node[lab,anchor=north] at (1.44,-2.5) {base window};
\draw[rulegray] (0.44,-2.5) -- (2.44,-2.5);
\end{tikzpicture}

\caption{What each gear is allowed to lean on. Notice the strong row needs a
\emph{window} of base cases: if the step reaches back $s$ positions, the
dominoes inside reach of no step must be toppled by hand.}
\label{fig:reachback}
\end{figure}

\begin{example}[Postage]\label{ex:postage}
Prove: every integer $n \ge 12$ can be written as $n = 3a + 7b$ for some
nonnegative integers $a, b$.
\par\smallskip
\textbf{Strategy.} Adding $3$ moves forward in steps of $3$, so the step will
reach back three positions: $P(n-2) \Rightarrow P(n+1)$. A reach of $3$ needs
a base window of $3$ consecutive cases --- check $12, 13, 14$ by hand.
\par\smallskip
\begin{solution}
Let $P(n)$ be ``$n = 3a + 7b$ for some $a, b \in \Z_{\ge 0}$''.
\emph{Base.} $12 = 3\cdot 4 + 7\cdot 0$, $13 = 3\cdot 2 + 7\cdot 1$,
$14 = 3\cdot 0 + 7\cdot 2$.
\emph{Step.} Let $n \ge 14$ and assume $P(m)$ for all $12 \le m \le n$. Then
$n+1 \ge 15$, so $n - 2 \ge 12$ and $n - 2 \le n$; by the I.H.,
$n - 2 = 3a + 7b$ for some $a, b \ge 0$. Then $n + 1 = 3(a+1) + 7b$.
\emph{Conclusion.} By strong induction, $P(n)$ holds for all $n \ge 12$.
\end{solution}
\end{example}

Note $n=11$ has no such representation --- $12$ is the honest base, not a
convenience. This is a baby Frobenius problem; the general theory waits for
Chapter~5's gcd machinery.

\begin{exercise}[medium]\label{exr:postage}
Prove every $n \ge 4$ is $2a + 5b$ for nonnegative $a,b$. How wide is your
base window?
\end{exercise}

\section{Well-ordering: induction in disguise}\label{sec:wop}
\index{well-ordering principle}

\begin{definition}[Well-ordering principle]
Every nonempty subset of $\N$ has a least element.
\end{definition}

WOP and induction are the same engine in different clothes. The translation is
the \emph{minimal counterexample}\index{minimal counterexample} template:
suppose some $n$ fails $P(n)$; take the \emph{smallest} failure $n_0$ (WOP
supplies it); then every smaller case holds, and the inductive step applied at
$n_0 - 1$ produces $P(n_0)$ --- contradiction. If you can prove the step but
induction bookkeeping annoys you, WOP is the same argument with the accounting
hidden.

\begin{theorem}\label{thm:primefactor}
Every integer $n \ge 2$ has a prime factor.
\end{theorem}

\begin{proof}
Suppose not, and let $n_0 \ge 2$ be the smallest counterexample (well-ordering).
Then $n_0$ is not prime --- otherwise it is its own prime factor --- so
$n_0 = ab$ with $2 \le a, b < n_0$. By minimality of $n_0$, the integer $a$ has
a prime factor $p$, and $p \mid a \mid n_0$ makes $p$ a prime factor of $n_0$:
contradiction.
\end{proof}

Compare with strong induction on the same claim: the ``smallest failure''
plays the role of $n+1$, and ``every smaller case'' is exactly the strong
hypothesis. Same proof, two costumes --- keep both, because some arguments
(descent, Section~\ref{sec:structure}'s fake proof) read more naturally as
minimal counterexamples.

\begin{figure}[ht]\centering
% Figure: well-ordering -- the smallest counterexample is where the argument lands.
% Source for Figure~\ref{fig:wop}. TikZ only; regenerable by compiling main.tex.
\begin{tikzpicture}[
  tick/.style={font=\footnotesize\sffamily,text=slate},
  ok/.style={circle,draw=rulegray,inner sep=1.2pt,minimum size=3.4mm},
  bad/.style={circle,draw=accent,fill=accent!20,inner sep=1.2pt,minimum size=3.4mm},
  lab/.style={font=\footnotesize\sffamily,text=slate,align=center}
]
\node[tick] at (-1.1,0.3) {claim fails at:};
\foreach \n/\x in {2/0,3/0.9,4/1.8,5/2.7,6/3.6,7/4.5,8/5.4,9/6.3} {
  \node[bad] at (\x,0.3) {\tiny \n};
}
\node[bad,very thick] (least) at (0,0.3) {\tiny 2};
\draw[->,very thick,accent] (least.north) -- ++(0,0.55)
  node[above,lab] {smallest failure $n_0$:\\ attack \emph{this} one};
\node[lab,align=center]{Every failing set has a least element (well-ordering).\\
Its minimality does all the work: every strictly smaller case is already fine.};
\end{tikzpicture}

\caption{Well-ordering finds the smallest failure. Notice: minimality is a
\emph{strong} induction hypothesis in disguise --- every smaller case is
already fine, which is exactly the ammunition the contradiction needs.}
\label{fig:wop}
\end{figure}

\section{Structural examples: a tiling, a recurrence, a fake proof}
\label{sec:structure}
\index{structural induction}

Induction is not only algebra. The object being inducted on can be a board, a
tree, a formula --- anything with a \emph{size} that decomposes into smaller
instances of itself.

\begin{example}[The tromino tiling]\label{ex:tromino}
An $L$-tromino covers three cells in an L shape. Prove: for every $n \ge 1$,
a $2^n \times 2^n$ board with any one cell removed can be tiled by L-trominos.
\par\smallskip
\textbf{Strategy.} Induct on $n$: the base $2 \times 2$ minus a cell \emph{is}
one tromino. For the step, the board must be made to look like four smaller
instances --- cut it into four $2^{n-1} \times 2^{n-1}$ quadrants. The hole
lives in one quadrant; place a single L-tromino on the three centre cells of
the other quadrants. Now \emph{every} quadrant is a $2^{n-1}$ board with one
cell spoken for --- the I.H. finishes each.
\par\smallskip
\begin{solution}
Let $P(n)$ be the claim for $2^n$ boards. \emph{Base.} $P(1)$: a $2 \times 2$
board missing one cell is exactly one L-tromino. \emph{Step.} Let $n \ge 2$
and assume $P(n-1)$: every $2^{n-1} \times 2^{n-1}$ board with one cell
removed is tileable. Given a $2^n \times 2^n$ board with a hole, quarter it.
The quadrant containing the hole is a $2^{n-1}$ board missing a cell; the
other three quadrants each lose one cell to the single centre tromino we place
(Figure~\ref{fig:tromino}, second panel). Apply the I.H. to all four.
\emph{Conclusion.} By induction, $P(n)$ holds for all $n \ge 1$.
\end{solution}
\end{example}

\begin{figure}[ht]\centering
% Figure: the tromino induction --- base, the quartering step, and the result.
% Source for Figure~\ref{fig:tromino}. TikZ only; regenerable by compiling main.tex.
\begin{tikzpicture}[
  cell/.style={draw=rulegray,minimum size=4.6mm,inner sep=0pt},
  hole/.style={fill=amberbg,draw=amber},
  ctr/.style={fill=accentsoft,draw=accent},
  lab/.style={font=\footnotesize\sffamily,text=slate},
]
% ---------- panel 1: base 2x2 ----------
\begin{scope}[xshift=0mm]
\foreach \p/\v in {0/0/{},0/1/B,1/0/B,1/1/B}
  \node[cell] at ({mod(\p,2)*0.46},{-int(\p/2)*0.46}) {\v};
\node[cell,hole] at (0,0) {};
\node[lab,anchor=north] at (0.46,-1.05) {base: $2\times2$};
\end{scope}
% ---------- panel 2: quartered 4x4 with centre tromino ----------
\begin{scope}[xshift=26mm]
\foreach \r in {0,...,3}
  \foreach \c in {0,...,3}
    \node[cell] at (\c*0.46,-\r*0.46) {};
\node[cell,hole] at (0,0) {};
% centre tromino (cells (1,2),(2,1),(2,2) in row,col)
\node[cell,ctr] at (2*0.46,-1*0.46) {};
\node[cell,ctr] at (1*0.46,-2*0.46) {};
\node[cell,ctr] at (2*0.46,-2*0.46) {};
% quarter boundaries
\draw[accent,thick] (0.92,0.23) -- (0.92,-1.15);
\draw[accent,thick] (-0.23,-0.92) -- (1.61,-0.92);
\node[lab,anchor=north,align=center]{cut in quarters,\\ one L at centre};
\end{scope}
% ---------- panel 3: subproblem view ----------
\begin{scope}[xshift=56mm]
\foreach \r in {0,...,1}
  \foreach \c in {0,...,1}
    \node[cell] at (\c*0.46,-\r*0.46) {};
\node[cell,hole] at (0,0) {};
\node[cell,ctr] at (1*0.46,-1*0.46) {};
\node[lab,anchor=north,align=center]{each quadrant:\\ a smaller hole};
\end{scope}
% ---------- panel 4: assembled 4x4 (real tiling) ----------
\begin{scope}[xshift=84mm]
\foreach \rc/\row in {0/{H,B,C,C},1/{B,B,A,C},2/{D,A,A,E},3/{D,D,E,E}}
  \foreach \c [count=\i] in \row
    \node[cell] at (\i*0.46-0.46,-\rc*0.46) {\c};
\node[cell,hole] at (0,0) {};
\node[lab,anchor=north] at (0.69,-1.3) {assembled: 5 trominos};
\end{scope}
\end{tikzpicture}

\caption{The induction at work. Notice what the step actually needs: one L
placed at the centre converts ``one quadrant has a hole'' into ``every quadrant
has a hole'' --- after which four copies of $P(n-1)$ finish the job. The last
panel is a real tiling produced by Listing~\ref{lst:tiling}'s algorithm.}
\label{fig:tromino}
\end{figure}

The proof is not just existence --- it is a construction recipe, and it runs:

\begin{checkpoint}[Predict before you run]
Listing~\ref{lst:tiling} tiles an $8 \times 8$ board with hole $(5,3)$. Before
running: predict the number of trominos ($(4^3 - 1)/3$ --- why is it an
integer?), and predict whether any cell can be covered twice.
\end{checkpoint}

\begin{lstlisting}[language=Python,caption={\texttt{verify\_tiling.py}: the
proof, compiled. Recursion \emph{is} the induction.},
label={lst:tiling}]
def tile(board, top, left, size, hr, hc, tid):
    if size == 2:                      # base: one L fills the non-hole cells
        ...
    h = size // 2
    q = (hr >= top+h)*2 + (hc >= left+h)      # quadrant holding the hole
    ...  # centre tromino covers one cell of each other quadrant
    for i, (qt, ql) in enumerate(origins):    # recurse into all four
        nhr, nhc = (hr, hc) if i == q else centers[i]
        tid = tile(board, qt, ql, h, nhr, nhc, tid)
\end{lstlisting}
The full script (in \texttt{code/verify\_tiling.py}) runs the construction on
$8\times 8$ and prints each tromino as a letter. Its real output:
\begin{lstlisting}[numbers=none]
board 8x8, hole at (5,3)
C C D D H H I I
C B B D H G G I
E B F F J J G K
E E F A A J K K
M M N N A R S S
M L N # R R Q S
O L L P T Q Q U
O O P P T T U U
trominos used: 21, cells covered: 63, expected: 21
\end{lstlisting}
Each letter is one tromino --- read the grid and you can \emph{see} the
quartering: the centre L is \texttt{A}, and each $4\times4$ quadrant repeats
the construction inside itself.

Induction is also how closed forms for recurrences get proved --- guess a
formula, then verify the step preserves it (Chapter~10 owns this; contest
sequences like Fibonacci identities are already fair game, e.g.\
AMC~12~2024~B~P18~\cite{aops-amc}).

Finally, the autopsy every student owes themselves once --- the most famous
broken induction ever written:

\begin{example}[Fake-proof autopsy: all horses one colour]\label{ex:horses}
``Theorem'': all horses are the same colour. ``Proof'': induct on the number
of horses. \emph{Base}: one horse has one colour. \emph{Step}: given $n+1$
horses, remove one --- the remaining $n$ share a colour by the I.H.; swap it
back and remove a different one --- again one colour; since both groups share
horses, all $n+1$ agree. Where does it break?
\par\smallskip
\textbf{Strategy.} Do not hunt for a subtle logical flaw --- hunt for the
smallest $n$ where the step's \emph{ingredients} stop existing. The argument
needs the two $(n)$-horse groups to \emph{share} a horse to transfer colour
through. When do they not?
\par\smallskip
\begin{solution}
At $n+1 = 2$: the groups $\{h_1\}$ and $\{h_2\}$ share nothing ---
Figure~\ref{fig:horses} marks the dead link. The step itself is genuinely
valid for every $n \ge 3$ (the shared middle horse transfers colour), but
validity alone never fires it: with no proof of $P(2)$, the link from
$P(1)$ to $P(2)$ is missing, so $P(2), P(3), P(4), \dots$ are all left
unestablished --- only $P(1)$ was ever proved. A single missing link kills
the entire chain from the break onward.
\end{solution}
\end{example}

\begin{figure}[ht]\centering
% Figure: the horses autopsy --- the step works for n>=3 and breaks at n=2.
% Source for Figure~\ref{fig:horses}. TikZ only; regenerable by compiling main.tex.
\begin{tikzpicture}[
  h/.style={circle,draw,minimum size=4.4mm,inner sep=0pt,font=\scriptsize},
  grp/.style={draw=accent,rounded corners=3pt,dashed,inner sep=2pt},
  lab/.style={font=\footnotesize\sffamily,text=slate},
  bad/.style={font=\footnotesize\sffamily,text=amber},
]
% ---- n=3: the step is fine ----
\node[lab] at (0,0.75) {$n=3$: step works};
\node[h,fill=accentsoft] (a) at (0,0) {};
\node[h,fill=accentsoft] (b) at (0.55,0) {};
\node[h] (c) at (1.1,0) {};
\draw[grp] (-0.28,0.28) rectangle (0.83,-0.28);
\draw[grp] (0.27,0.45) rectangle (1.38,-0.11);
\node[lab,anchor=west] at (1.55,0.35) {$h_1,h_2$ same colour};
\node[lab,anchor=west] at (1.55,-0.1) {$h_2,h_3$ same colour};
\node[lab,anchor=west] at (1.55,-0.55) {shared horse $h_2$ $\Rightarrow$ all equal};
% ---- n=2: the link is broken ----
\node[lab] at (0,-1.9) {$n=2$: link broken};
\node[h] (x) at (0,-2.65) {};
\node[h] (y) at (0.55,-2.65) {};
\draw[grp] (-0.28,-2.37) rectangle (0.28,-2.93);
\draw[grp] (0.27,-2.2) rectangle (0.83,-2.76);
\node[bad,anchor=west] at (1.55,-2.4) {no shared horse};
\node[bad,anchor=west] at (1.55,-2.85) {colour cannot transfer};
\draw[amber,thick] (0.415,-2.45) -- (0.415,-2.85);
\node[amber,font=\normalsize] at (0.415,-3.1) {\textbf{\texttimes}};
\end{tikzpicture}

\caption{Why the horses proof dies at $n=2$. Notice the colour-transfer
mechanism is real for $n \ge 3$ --- the fake proof fails on the \emph{reach}
of the step, not on its logic. This is the commonest way inductions break in
the wild.}
\label{fig:horses}
\end{figure}

\begin{exercise}[medium]\label{exr:fake}
A ``proof'' claims every $n \ge 0$ satisfies $n = 0$, because $P(0)$ holds and
$P(n) \Rightarrow P(n+1)$ is proved by ``adding nothing to both sides of
$n = 0$''. Identify the actual failure.
\end{exercise}

\section{Choosing what to induct on}\label{sec:choosevar}
\index{induction!choosing the variable}

The template asks one question most students skip: \emph{what is $n$?}
The statement's variable is the default, but any quantity that measures the
problem's size works --- board dimension, set cardinality, string length,
number of operators. Rule of thumb: induct on the quantity the problem
\emph{decomposes along}. A $2^n$ board splits along $n$; a graph along $|V|$;
a sum along its term count.

A subtler choice: the \emph{claim itself}. Sometimes $P(n)$ is too weak to
carry the step, and the fix is \emph{inductive loading} --- prove a
\emph{stronger} statement so the I.H. gives you more. Classic instance: bound
$\sum_{k=1}^{n} 1/k^2 < 2$ directly and the step stalls (the surplus shrinks);
prove the stronger $\sum_{k=1}^{n} 1/k^2 \le 2 - 1/n$ and the step goes
through cleanly. Chapter~13 uses this move for inequalities.

\begin{exercise}[easy]\label{exr:choose}
For each claim, name what you would induct on: (a) every full binary tree with
$\ell$ leaves has $2\ell - 1$ vertices; (b) every $n \ge 60$ is a sum of
$6$'s, $9$'s and $20$'s; (c) a formula with $m$ operators has at most $m+1$
variables.
\end{exercise}

\section{Chapter summary}
\begin{table}[ht]\centering
\begin{tabularx}{\linewidth}{@{}lYYY@{}}
\toprule
\tablehead\thead{Situation} & \thead{Gear} & \thead{Base obligation} & \thead{Step obligation}\\
\midrule
standard $\forall n \ge b$ & induction & $P(b)$ & $P(n) \Rightarrow P(n+1)$\\
step reaches back $s$ & strong induction & $s$ consecutive bases & all of $P(b..n) \Rightarrow P(n+1)$\\
``no smallest bad case'' style & well-ordering & --- & smallest failure contradicts itself\\
algebraic sum/formula & peel last term & check by hand & split so the I.H. appears\\
inequality (eventual) & squeeze & the \emph{first} true $n$ & name every $\ge$ link\\
structure that decomposes & structural & smallest instance & reassemble smaller copies\\
claim too weak for the step & strengthen $P(n)$ & stronger base & stronger I.H. carries it\\
\bottomrule
\end{tabularx}
\caption{The induction menu on one page. The right column is where each
variant actually breaks in student work --- check it before submitting.}
\label{tab:induction}
\end{table}

\section{Common mistakes}
\index{common mistakes}
\begin{itemize}
\item \textbf{No base case --- or the wrong window.} The horses proof dies
exactly here: a step valid only from $n \ge 3$ needs $P(1)$ \emph{and} $P(2)$
checked by hand. Always ask: how far back does my step reach? That distance
is your base window.
\item \textbf{Using the conclusion inside the step.} ``Assume $P(n+1)$\dots''
or quietly writing $\sum_{k=1}^{n+1} = \frac{(n+1)(n+2)}{2}$ mid-derivation is
assuming what you must prove. The I.H. covers $P(n)$ (or earlier), never
$P(n+1)$.
\item \textbf{Inducting on the wrong variable.} On a two-parameter claim
($2^n \times 2^n$ boards, $n \times m$ grids), inducting on the wrong
coordinate gives a step that cannot decompose the object. Induct on the
quantity the structure splits along (Section~\ref{sec:choosevar}).
\end{itemize}

\section{Additional exercises}
\begin{exercise}[easy]
Prove by induction: $\sum_{k=1}^{n} (2k - 1) = n^2$. (What does the picture
of $n^2$ as nested Ls say?)
\end{exercise}
\begin{exercise}[easy]
Prove by induction: $\sum_{k=1}^{n} k \cdot k! = (n+1)! - 1$.
\end{exercise}
\begin{exercise}[easy]
Prove: $3 \mid (4^n - 1)$ for all $n \ge 0$. (This is why $21 = (4^3-1)/3$
trominos was guaranteed to be an integer.)
\end{exercise}
\begin{exercise}[easy]
Each of these inductions is broken --- name the failure type from Common
mistakes: (a) a proof that all integers are equal (``true for $n$ implies
true for $n+1$ by symmetry''); (b) a proof of $\sum k^3 = n^4$ whose step
``assumes the formula for $n+1$ and verifies it''.
\end{exercise}
\begin{exercise}[medium]
Prove: $n^3 - n$ is divisible by $6$ for all $n \ge 0$, once by induction and
once by the factorisation $n^3 - n = (n-1)n(n+1)$.
\end{exercise}
\begin{exercise}[medium]
Prove by induction: $F_1 + F_2 + \cdots + F_n = F_{n+2} - 1$, where
$F_1 = F_2 = 1$. (Fibonacci sums are AMC staples --- see
AMC~12~2024~B~P18~\cite{aops-amc}.)
\end{exercise}
\begin{exercise}[medium]
Prove: every $n \ge 1$ has a binary representation $n = \sum 2^{a_i}$ with
distinct exponents, by strong induction. (Halve $n$ and read the parity; the
$n=1$ base is $2^0$.)
\end{exercise}
\begin{exercise}[medium]
Prove by minimal counterexample: every rational number has a lowest-terms
representative $p/q$ with $\gcd(p,q) = 1$. (Take the smallest denominator.)
\end{exercise}
\begin{exercise}[hard]
Prove that for every $n \ge 1$, a $2^n \times 2^n$ board missing \emph{two}
cells can never be tiled by L-trominos. (Area: $(4^n - 2)/3$ must be an
integer --- prove $4^n \equiv 1 \pmod 3$ by induction, so the count is never
one.)
\end{exercise}
\begin{exercise}[hard]
Extend \texttt{code/verify\_tiling.py} so the hole can be \emph{any} cell and
the board any $2^k$ (try $16 \times 16$). Then prove the count invariant:
the printed number of trominos is always $(4^n - 1)/3$, and prove by
induction that $3 \mid 4^n - 1$ so the count is always an integer.
\end{exercise}
\begin{exercise}[hard]
Strengthened induction. Prove directly that $\sum_{k=1}^{n} 1/k^2 < 2$
stalls, then prove the stronger $\sum_{k=1}^{n} 1/k^2 \le 2 - 1/n$ by
induction. Explain in one sentence why a stronger claim can be easier.
\end{exercise}
\begin{exercise}[easy]
Prove by induction: $1 + 3 + 5 + \cdots + (2n-1) = n^2$. (Draw the
staircase of Figure~\ref{fig:stairsum} first.)
\end{exercise}
\begin{exercise}[medium]
Prove by induction that $6 \mid n^3 - n$ for every $n \ge 0$. Which
property of $n^3 - n$ does the step use at the end?
\end{exercise}
\begin{exercise}[medium]
Prove by \emph{strong} induction: every integer $n \ge 12$ is a sum of
$4$'s and $5$'s. State exactly which smaller cases the step reaches back
to (Figure~\ref{fig:wop} is the right picture).
\end{exercise}
\begin{exercise}[hard]
Find the flaw: ``all horses are the same colour'' is proved by induction
--- $n=1$ trivial; given $n$ same-coloured horses, in any set of $n{+}1$
the first $n$ share a colour and the last $n$ share a colour, so all
$n{+}1$ do. Name the exact value of $n$ where the step fails and why
the overlap argument collapses there. (Compare Figure~\ref{fig:horses}.)
\end{exercise}

\section*{Further reading}
\begin{itemize}
\item E.~Lehman, F.~T.~Leighton, A.~R.~Meyer, \emph{Mathematics for Computer
Science}, Ch.~5 ---
\url{https://courses.csail.mit.edu/6.042/spring18/mcs.pdf} --- free PDF;
ordinary, strong and structural induction plus the well-ordering
equivalence this chapter quotes.
\item R.~Hammack, \emph{Book of Proof}, 3rd ed., Ch.~10 ---
\url{https://richardhammack.github.io/BookOfProof/} --- free PDF; the
clearest gallery of induction proofs, including the inequality examples
this chapter borrows.
\item MIT OpenCourseWare, \emph{6.042J Mathematics for Computer Science}
(Fall 2010), Lectures 2--3 ---
\url{https://ocw.mit.edu/6-042JF10} --- the induction lectures;
the ``bogus proofs'' tradition our autopsy section continues.
\item Stanford CS103, \emph{Mathematical Foundations of Computing} ---
\url{https://web.stanford.edu/class/cs103/} --- free reader; its induction
guide is the best drill on base-window discipline.
\item The Open Logic Project ---
\url{https://openlogicproject.org/} --- free texts; for readers who want
induction as a formal deduction rule rather than a template.
\item Cut-the-Knot, \emph{Induction} pages ---
\url{https://www.cut-the-knot.org/} --- the classic fallacies collection;
the fake-proof genre of \S\ref{sec:structure} lives here.
\end{itemize}

\begin{rulebox}[Watch instead]
\begin{itemize}
\item \textbf{MIT OpenCourseWare --- ``MIT 6.042J Mathematics for Computer
Science''} (playlist; induction at Lectures 2--3),
\url{https://www.youtube.com/playlist?list=PLB7540DEDD482705B} ---
covers \S\ref{sec:domino}--\ref{sec:strong}: the induction axiom,
worked sums, and strong induction. (Verified 2026-10-11.)
\item \textbf{SawFin Mathematics --- ``Discrete Math II -- 5.2.1 Proof by
Strong Induction''}, 16:21,
\url{https://www.youtube.com/watch?v=-vJhahLdmLI} --- covers
\S\ref{sec:strong}: when the single-predecessor step is not enough.
(Verified 2026-10-11.)
\end{itemize}
\end{rulebox}

\chapternotes
The domino template, the base-window discipline, and the strong-vs-ordinary
reformulation follow MIT 6.042J's induction and well-ordering chapters, which
also supplied the WOP$\,\Leftrightarrow\,$induction equivalence and the
``bogus proofs'' tradition our autopsy section continues~\cite{mcs}.
Velleman's induction chapter organises proofs by the same
strategy-first pattern this book uses~\cite{velleman}, and Hammack's chapter 10
supplied the inequality-by-induction examples~\cite{hammack}. The tromino
theorem is the classic first structural induction (a Cut-the-Knot staple;
the constructive proof here is the standard
one~\cite{cuttheknot}); Engel's induction chapter catalogues the
choosing-the-variable and strengthen-the-claim moves of
Section~\ref{sec:choosevar}~\cite{engel}. Contest sequence problems named in
this chapter (AMC~12~2024~A~P21, AMC~12~2024~B~P18, AIME~2021~I~P10,
AIME~2025~II~P13, AMC~12~2023~A~P20, AMC~12~2023~A~P25) live on the AoPS
archive~\cite{aops-amc} --- they are inductions wearing arithmetic clothes.
